% ============================================================================
% CHAPTER 7: Getting Good Information from People
% Why every label is an argument — and why that argument matters
% ============================================================================
% Source: /markdown/ch7/Ch7_final.md
% ============================================================================

\begin{chapterquote}{The premise of Chapter 7}
Human labels are the foundation of supervised learning---but labels are not neutral. Every annotation is an act of interpretation, shaped by who is labeling, what instructions they received, and what world they live in.
\end{chapterquote}

% ============================================================================
\section{The Photograph That Changed How We See Data}
\label{sec:ch7-imagenet}
% ============================================================================

In 2010, a Stanford researcher named Fei-Fei Li stood in front of an audience at a TED conference in Long Beach and explained what she had spent four years building: ImageNet. The project, completed using Amazon Mechanical Turk, had produced 15 million photographs organized into 22,000 categories. It was the largest labeled image dataset ever assembled, and it would become the training ground for the deep learning revolution that followed \autocite{deng2009imagenet}.

By 2019, a series of studies had identified something troubling in the data. The category ``cheerleader'' contained photographs that were overwhelmingly of women. The category ``guard'' was overwhelmingly men. The annotations were accurate in a narrow sense---the photographs did show what they were labeled as---but the dataset had encoded cultural stereotypes so thoroughly that models trained on it would confidently predict ``nurse'' for women and ``doctor'' for men when given ambiguous inputs \autocite{yang2018fairer}.

The labels were not wrong. They were incomplete descriptions of a biased sample, faithfully transcribed into a format that the model interpreted as ground truth.

The ImageNet case illustrates something that the whole field of machine learning spent the 2010s slowly learning: \term{labels are not neutral}. Every label is an act of interpretation, shaped by who is doing the labeling, what instructions they were given, what cultural context they bring, and how tired they are. The data annotation problem is not a logistical problem that can be solved by deploying enough workers. It is an epistemological problem---a problem about how we know what we know, and what it means to reduce the complexity of a photograph, a sentence, or a human interaction to a single numerical label.

% ============================================================================
\section{The Instruction Problem}
\label{sec:ch7-instructions}
% ============================================================================

Imagine you are about to annotate a set of online comments for toxicity. You have been given a single instruction: ``Label each comment as TOXIC or NOT TOXIC.''

You see this comment: \textit{``I can't believe you people still think this is acceptable. You're all completely delusional.''}

Is it toxic? Probably yes, you think.

Now imagine the same comment, but a second annotator is working with a slightly expanded instruction: ``Label each comment as TOXIC or NOT TOXIC. Toxicity means explicit slurs, direct threats, or content that would make a member of a protected group feel physically unsafe.''

The second annotator looks at the same comment. Contemptuous? Yes. An explicit slur, a direct threat, or content likely to make a protected group feel physically unsafe? Not exactly. They label it NOT TOXIC.

Same comment. Different instructions. Different labels.

This is what annotation researchers call the \term{instruction problem}, and it is pervasive. A study by Cer and colleagues at Google examined how different framings of the same toxicity task produced dramatically different label distributions even when the underlying comments were identical. The framing ``Is this comment hateful?'' produced more toxic labels than the framing ``Would a reasonable person find this comment toxic?'' The second framing invokes a hypothetical external standard that moderates individual variation---it asks the annotator to model a composite person rather than report their own reaction.

The instruction problem compounds when guidelines are long. Professional annotation companies routinely produce annotation guidelines that run to 50, 80, even 120 pages for complex labeling tasks. Annotators working from a 100-page guideline develop a mental model of the guideline that is partly shared and partly individual. The shared part produces consistent labels. The individual part produces the kind of systematic variation that shows up in low inter-annotator agreement scores.

% ============================================================================
\section{Measuring Agreement: Cohen's Kappa}
\label{sec:ch7-kappa}
% ============================================================================

When two annotators label the same set of items, their agreement is measurable. The simplest measure is raw agreement---the proportion of items they label the same way. But raw agreement is inflated by chance agreement: if 90\% of items are NOT TOXIC and both annotators are just labeling everything NOT TOXIC, they will agree 81\% of the time by chance, while making no genuine classification decisions.

The standard correction is \term{Cohen's kappa}, developed by psychologist Jacob Cohen in 1960:

\[
\kappa = \frac{P_o - P_e}{1 - P_e}
\]

where $P_o$ is the observed agreement and $P_e$ is the expected agreement by chance. A kappa of 1.0 means perfect agreement; 0 means agreement no better than chance; negative values mean systematic disagreement.

\begin{keyconcept}[The Landis-Koch Interpretation Scale]
\begin{itemize}
    \item $\kappa < 0.20$: Slight agreement
    \item $0.20 \leq \kappa < 0.40$: Fair agreement
    \item $0.40 \leq \kappa < 0.60$: Moderate agreement
    \item $0.60 \leq \kappa < 0.80$: Substantial agreement
    \item $\kappa \geq 0.80$: Almost perfect agreement
\end{itemize}
These thresholds are conventional, not theoretically grounded. Always interpret kappa in the context of the specific task.
\end{keyconcept}

When the toxicity detection field began systematically measuring annotator agreement, the results were sobering. A comprehensive study by Zampieri and colleagues in 2019 found that annotation kappa for English-language toxicity tasks varied between 0.45 and 0.72 depending on the task framing and annotator pool \autocite{zampieri2019predicting}. The lower end---``moderate agreement''---meant that two annotators looking at the same comment would disagree about a quarter of the time. A model trained on such data is not learning to detect toxicity; it is learning to predict the label distribution of a specific group of annotators working under specific conditions.

% ============================================================================
\section{The Annotator Is Not a Recording Device}
\label{sec:ch7-annotator-bias}
% ============================================================================

There is a useful fantasy in machine learning that annotators are passive transcribers---that they observe some objective feature of the data and record it, and that the label accurately reflects something real in the world. The fantasy is comfortable because it preserves the idea that labels are ground truth.

It is false.

An annotator is a person with a history, a worldview, a cultural context, and a present state that varies by hour of day, level of fatigue, number of items reviewed, and recent sequence of items seen.

\paragraph{Sequence effects.} An annotator who has just reviewed 20 highly toxic comments will have a different calibration threshold for the next borderline comment than an annotator who has just reviewed 20 benign ones. Studies by Callison-Burch at the University of Pennsylvania, examining machine translation quality annotation, showed that annotators systematically rated the same translations differently depending on the quality of the surrounding translations in their queue \autocite{callisonburch2010creating}. The translation had not changed. The annotator's reference point had.

\paragraph{Demographic effects.} Research by Sap and colleagues at Carnegie Mellon found that Black annotators and white annotators differed substantially in their labeling of tweets containing African American Vernacular English (AAVE). Content in AAVE---including language that was neutral or positive in context---was more likely to be labeled as toxic when the annotator pool was predominantly white \autocite{sap2019risk}. The consequence is that a model trained on annotations from a predominantly white annotator pool will have a systematically higher false toxicity rate for AAVE.

\begin{cautionbox}[The Demographics-Bias Connection]
This is not a problem that better annotation guidelines can fully fix, because it is embedded in who the annotators are and what cultural context they bring to the task. Diversifying the annotator pool is a practical necessity, not merely a fairness gesture.
\end{cautionbox}

% ============================================================================
\section{The Economics of Getting Labels}
\label{sec:ch7-economics}
% ============================================================================

The ImageNet dataset was completed with the help of Amazon Mechanical Turk, a crowdsourcing platform that allows requesters to post small tasks and have them completed by a globally distributed workforce for small payments. A typical toxicity labeling task on Mechanical Turk pays \$0.02 to \$0.05 per item.

The economic logic of crowdsourcing annotation is compelling. A dataset that would cost hundreds of thousands of dollars to label through professional linguists can be labeled for tens of thousands of dollars through Mechanical Turk, at higher speed.

The quality control problem is also real. Without oversight, crowdsourced annotation produces noise: workers who click randomly, who misread instructions, whose cultural context does not match the task requirements. The standard approach---multiple annotators per item plus majority voting---reduces individual noise. But majority voting on disagreement can enshrine the biases of the majority as ground truth.

The distinction between \term{noise disagreement} and \term{genuine ambiguity} matters enormously. If two annotators disagree because one was not paying attention, the disagreement is noise and should be resolved by majority vote. If two annotators disagree because the item is genuinely ambiguous, the disagreement is information. Collapsing it to a single label loses the information.

\begin{examplebox}[Label Distribution Training]
Some modern machine learning practice has begun to work with label distributions rather than single labels---training on the full distribution of annotator judgments rather than the majority vote. Research by Davani and colleagues demonstrated that models trained on distributions were better calibrated: they expressed more uncertainty on genuinely ambiguous items, and more confidence on items where annotators had agreed. This is a more honest representation of what the data actually contains~\autocite{davani2022dealing}.
\end{examplebox}

% ============================================================================
\section{Framing Effects: How the Question Changes the Answer}
\label{sec:ch7-framing}
% ============================================================================

Framing effects describe the phenomenon that how a question is posed shapes the answer, independently of what the answer ``truly'' is. Amos Tversky and Daniel Kahneman showed that a medical treatment described as having a ``90\% survival rate'' was viewed more favorably than the same treatment described as having a ``10\% mortality rate.'' Same statistic. Very different responses \autocite{tversky1981framing}.

The same dynamic appears in annotation. Consider the difference between:

\begin{center}
\textit{Version A:} ``Rate each response on a scale of 1--5 for helpfulness.''
\end{center}

\begin{center}
\textit{Version B:} ``Consider: would a patient receiving this information be better equipped to make decisions about their health? Rate 1--5.''
\end{center}

Version B anchors the annotator on a concrete downstream beneficiary and asks them to simulate that person's experience. This produces more consistent labels across annotators and labels that better predict actual user satisfaction. The question is not whether Version B is more ``correct.'' It is that Version B more precisely specifies what the annotation is supposed to measure.

The most consequential framing decisions in annotation are often invisible because they are embedded in what is \emph{not} included in the instructions. Medical image annotation typically asks annotators to identify the presence or absence of a finding, not to classify its clinical significance. A model trained on such annotations learns to detect findings, not to assess their relevance to patient management---a critical difference when the model is deployed to assist diagnosis.

% ============================================================================
\section{Calibration Sessions and the Iterative Guideline}
\label{sec:ch7-calibration}
% ============================================================================

The annotation community has developed practices to address the instruction problem, the most important of which is the \term{calibration session}.

In a calibration session, a group of annotators reviews and discusses a set of carefully selected items together. The items are chosen to cover the range of cases that arise in the task, including clear cases, borderline cases, and cases that expose ambiguities in the annotation guidelines. Annotators record their labels independently, then compare and discuss disagreements.

The output is twofold: an updated annotation guideline that addresses the exposed ambiguities, and a shared mental model that enables more consistent application going forward. Companies like Scale AI and Appen conduct formal calibration sessions at the start of annotation projects and regularly thereafter, particularly when the task involves subjective judgment or cultural nuance.

OpenAI's work on RLHF (Reinforcement Learning from Human Feedback) is explicit about this iterative process. The human feedback that shapes RLHF training is not collected once and used permanently; it is continuously re-evaluated, guidelines are updated, and new annotation rounds are conducted to capture changing standards \autocite{ouyang2022training}. The annotation is not a solved problem but an ongoing conversation between the model's behavior and human preferences.

% ============================================================================
\section{What Low Agreement Tells You}
\label{sec:ch7-low-kappa}
% ============================================================================

When inter-annotator agreement on a task is low---say, kappa below 0.40---the conventional response is to interpret this as a quality problem. The solution is to revise the guidelines and improve the training.

This interpretation is sometimes correct. But low agreement can also be a signal about the \emph{nature of the task itself}. Some annotation tasks ask annotators to make distinctions that are genuinely difficult, genuinely subjective, or genuinely contested. A task asking annotators to distinguish between ``mildly offensive'' and ``moderately offensive'' is asking for a discrimination that may not be reliably present in human judgment, not because the annotators are failing but because the categories do not map onto stable psychological reality.

When low agreement persists despite careful guidelines, calibrated training, and expert annotators, the right response is often not to force agreement but to interrogate the task. Is the distinction the annotation is trying to capture actually important for the downstream use case? Could the model perform adequately with coarser categories where agreement is higher? Is the genuine ambiguity in the data actually carrying information that should be preserved?

\begin{keyconcept}[A Low Kappa Score Is Telling You Something Real]
In the most honest framing: a low kappa score is telling you something real about the data, about the task, and about the limits of reducing human judgment to a scalar label. The answer to that signal is not always to try harder. Sometimes it is to listen.
\end{keyconcept}

% ============================================================================

\begin{trythis}
Take any three subjective opinions you hold---about a film, a piece of writing, a public figure's behavior. For each one, write down the implicit criteria you are using to form the opinion. Now try to write those criteria as annotation guidelines, clearly enough that a stranger could apply them reliably. Notice how quickly ``good'' requires unpacking, how ``harmful'' requires definition, and how much of your judgment is tacit knowledge that resists specification. This is the annotation problem, and it does not have a clean solution.
\end{trythis}

% ============================================================================
\section*{Chapter Summary}
\addcontentsline{toc}{section}{Chapter Summary}
% ============================================================================

\begin{keyconcept}[Key Concepts]
\begin{itemize}
    \item Labels are not neutral---every annotation is an act of interpretation shaped by who is labeling, what instructions they received, and what cultural context they bring
    \item The instruction problem: even small differences in guideline wording produce systematically different label distributions
    \item \term{Cohen's kappa} measures inter-annotator agreement corrected for chance; low kappa can signal poor guidelines, genuine task ambiguity, or both
    \item Framing effects show that the question shape changes the answer, independently of what is being classified
    \item The annotator is not a recording device---sequence effects, demographic background, and fatigue all systematically shape labels
\end{itemize}
\end{keyconcept}

\subsection*{Key Examples}

\begin{description}
    \item[ImageNet (2010--2019)] Cultural stereotypes baked into labels by an annotation process that faithfully transcribed annotator worldview as ground truth
    \item[AAVE toxicity labeling] Demographically homogeneous annotator pools produce systematically biased labels for minority language registers
    \item[RLHF human feedback (OpenAI)] Iterative guideline revision as a necessary feature of annotation for \ai training
    \item[Scale AI and Appen calibration sessions] The standard industry practice for reducing annotation inconsistency
    \item[Label distribution training] Training on the full annotation distribution rather than majority-vote labels produces better-calibrated models
\end{description}

\bigskip
\begin{center}
\rule{0.5\textwidth}{0.4pt}
\end{center}
\bigskip

\noindent\textit{In the next chapter, we zoom out from individual interactions---interfaces, labels, human judgments---and look at the \hitl pipeline as a whole: the data flows, the engineering challenges, and what happens when a system that worked beautifully in the lab breaks down in production.}
